\documentclass[12pt,a4paper]{article}
\usepackage[utf8]{inputenc}
\usepackage[T1]{fontenc}
\usepackage[english]{babel}
\usepackage{amsmath}
\usepackage{geometry}
\usepackage[hidelinks]{hyperref}
\usepackage{enumitem}
\geometry{top=2.5cm,bottom=2.5cm,left=2.7cm,right=2.7cm}
\setlist[itemize]{leftmargin=*,topsep=2pt,itemsep=2pt,parsep=0pt}
\setlength{\parindent}{0pt}
\setlength{\parskip}{5pt}
\begin{document}
\begin{center}
{\Large\bfseries European Semiconductor Financial Benchmarking}\\[0.65em]
{\normalsize Salvatore Megna}\\[0.25em]
{\small Data Science / Statistical Research}
\end{center}
\vspace{1.1em}
\par
\vspace{1.15em}
\noindent{\bfseries Research question}\par
How do six European semiconductor companies compare in growth, profitability, innovation spending and balance-sheet strength?
\par
\par
\vspace{1.15em}
\noindent{\bfseries Data and methods}\par
Financial-statement panel for ASML, ASM International, Besi, AIXTRON, VAT Group and Technoprobe. Main sample: 2021--2025, 30 company-years, with 2020 as an auxiliary base year.

Six financial ratios, annual median benchmarks, IQR screening and PCA on five standardized indicators; original workbook and source records retained.
\par
\par
\vspace{1.15em}
\noindent{\bfseries Selected results}\par
\begin{itemize}
\item Median revenue growth: 31.70\% (2021), -5.33\% (2023) and 11.07\% (2025).
\item PCA explains 46.17\% on PC1 and 26.47\% on PC2 (72.65\% combined).
\item IQR screening highlights Besi 2021 and VAT Group 2022, without removing them.
\end{itemize}
\par
\par
\vspace{1.15em}
\noindent{\bfseries Interpretation and limitations}\par
Six firms do not constitute a representative industry sample. Different reporting policies and currencies matter, and the PCA describes correlations rather than causes.
\par
\par
\vspace{1.15em}
\noindent{\bfseries Reproduce}\par
\texttt{Rscript scripts/run\_all.R}
\par
\vfill
\noindent\small\href{https://github.com/Salvatore-Megna/european-semiconductor-benchmarking}{Repository and full documentation}
\end{document}
